\documentclass[10pt,a4paper]{amsart}
\usepackage[margin=19mm]{geometry}
\usepackage{amsmath,amssymb,mathtools}
\usepackage[scr=rsfs,cal=euler]{mathalfa}
\usepackage{xfrac}
\usepackage[unicode,hidelinks,bookmarks=false]{hyperref}
\newcommand{\ie}{i.e.\ }
\newcommand{\cf}{cf.\ }
\newcommand{\resp}{respectively\ }
\newcommand{\Subsection}{\subsection}
\providecommand{\nobreakdash}{\nobreak\hbox{-}}
\setlength{\emergencystretch}{2em}
\multlinegap0pt
\usepackage{algorithm}\usepackage{algpseudocode}
\usepackage{amssymb}
\newtheorem{proposition}{Proposition}
\newtheorem{lemma}{Lemma}
\newcommand{\C}{\mathbb{C}}
\newcommand{\Normal}{\mathcal{N}}
\newcommand{\R}{\mathbb{R}}
\newcommand{\distributions}{\mathcal{M}_1} %
\newcommand{\hs}[1]{\mathcal{#1}} %
\newcommand{\stp}[2]{{\overset{(#2)}{#1}}}
\newcommand{\stpx}[1]{{(#1)}}
\title{Statistical Learning Theory for Distributional Classification}
\author{Christian Fiedler}
\date{Source: arXiv:2601.14818v1 (CC BY 4.0)}
\begin{document}
\maketitle

\noindent\emph{Source and licence.} Statistical Learning Theory for Distributional Classification. Version arXiv:2601.14818v1 (CC BY 4.0), distributed by arXiv under the Creative Commons Attribution 4.0 International licence, http://creativecommons.org/licenses/by/4.0/, as recorded in the arXiv licence metadata for this exact version and shown on the abstract page https://arxiv.org/abs/2601.14818v1. Full licence notice: \texttt{licenses/arxiv-2601.14818-CC-BY-4.0.txt}.\par\smallskip

\noindent\textit{Editorial scope.} Counted unit: the article's lemma lemma (source0.tex lines 1629--1632). The article's complete original proof of it is delivered with it (source0.tex lines 1633--1684, one \texttt{proof} environment of 52 lines). No mathematics is written, repaired or re-derived: the statement and the proof are the article's own lines, in the article's own notation, and no retained line is substituted or re-flowed.  Retained uncounted as the source-local context the proof needs: lines 1567--1574.  Nothing was omitted from any retained span, so the package carries no omission record.  No retained line contains a citation, so the wrapper carries no bibliography and no bibliography entry.  No retained line contains a figure environment, an image-inclusion command or drawing code, so no image is delivered and the package declares no asset dependency.  Editorial formatting only, in addition: an AMS \texttt{amsart} preamble that restates the article's own declarations and macros used by the retained lines, namely the environments \texttt{proposition}, \texttt{lemma} on separate counters, together with the standard packages those lines need.  The retained environments print in this document as Proposition 1, Lemma 1 in order of appearance, whereas the article prints the same items with the numbering of its own full document.  The complete native source of the article as distributed in the arXiv e-print for this exact version is bundled unchanged as \texttt{sources/193112/source0.tex}.
\begin{proposition} \label{prop:gaussian:existenceUniquenessFourier}
For all $a\in\hs{H}$ and $Q\in L_1^+(\hs{H})$ there exists exactly one Borel probability measure $\mu\in\distributions(\hs{H})$ with
\begin{equation}
    \int_\hs{H} \exp(i\langle h,x\rangle) \mathrm{d}\mu(x) = \exp(i\langle a, h\rangle_\hs{H})\exp\left(-\frac12 \langle Qh, h\rangle_\hs{H}\right)
    \quad \forall h\in\hs{H}
\end{equation}
This probability measure has mean $a$, covariance operator $Q$, and we denote it by $\Normal(a,Q)$.
\end{proposition}

\begin{lemma}
Let $\hs{H}$ be a separable real Hilbert space, $Q\in L_1^+(\hs{H})$ with $\ker(Q)=\{0\}$, and $\lambda\in\R$.
The map $\hs{H}\ni h \mapsto \exp(i W_h(\cdot)) \in L^2(\hs{H},\Normal(0,Q);\C)$ is well-defined and continuous.
\end{lemma}

\begin{proof}
Recall that since $\ker(Q)=\{0\}$, $Q^\frac12(\hs{H})$ is dense in $\hs{H}$.
Let $\lambda\in\R$ and $h\in Q^{\frac12}(\hs{H})$. We have
\begin{align*}
    \exp\left(-\frac{\lambda^2}{2}\|h\|_\hs{H}^2\right)
    & = 1 \cdot \exp\left(-\frac{1}{2}\langle \lambda Q^{\frac12}Q^{-\frac12}h, \lambda Q^{\frac12}Q^{-\frac12}h \rangle_\hs{H}\right)  \\
    & \stp{=}{1} \exp(i \langle 0, h\rangle_\hs{H}) 
        \exp\left(-\frac{1}{2}\langle \lambda Q \left(Q^{-\frac12}h\right), \lambda \left(Q^{-\frac12}h\right) \rangle_\hs{H}\right) \\
    & \stp{=}{2} \int_\hs{H} \exp(i\langle  \lambda Q^{-\frac12}h, x \rangle_\hs{H})\mathrm{d}\Normal(x\mid 0, Q) \\
    & = \int_\hs{H} \exp(i\lambda \langle Q^{-\frac12}h, x \rangle_\hs{H})\mathrm{d}\Normal(x\mid 0, Q) \\
    & \stp{=}{3} \int_\hs{H} \exp(i\lambda W_h(x))\mathrm{d}\Normal(x\mid 0, Q),
\end{align*}
where we used for \stpx{1} the fact that $Q^{\frac12}$ is self-adjoint and $e^0=1$,
in \stpx{2} we used the characterization of the Gaussian measure (via the Fourier transform) from Proposition \ref{prop:gaussian:existenceUniquenessFourier},
and in \stpx{3} the definition of $W_h(x)$ (recall that $h\in Q^{\frac12}(\hs{H})$ by assumption).
Note that since $|e^{is}|=1$ for all $s\in\R$, $\exp(i \lambda W_h(\cdot))\in L^2(\hs{H},\Normal(0,Q);\C)$.

Next, we show that for all $\lambda\in\R$, the map $Q^{\frac12}(\hs{H})\ni h \mapsto \exp(i \lambda W_h) \in L^2(\hs{H},\Normal(0,Q);\C)$ is continuous. 
Since $Q^{\frac12}(\hs{H})$ is a vector space, and $h \mapsto W_h$ is linear, it is enough to show this for $\lambda=1$.
Let $h_1,h_2\in Q^{\frac12}(\hs{H})$, then
\begin{align*}
    & \|\exp(iW_{h_1}(\cdot)) - \exp(iW_{h_2}(\cdot))\|_{L^2(\hs{H},\Normal(0,Q);\C)}^2
    = \int_\hs{H} \left|\exp(iW_{h_1}(x)) - \exp(iW_{h_2}(x))\right|^2\mathrm{d}\Normal(x\mid 0,Q) \\
    & \hspace{0.5cm} = \int_\hs{H} \left(\exp(iW_{h_1}(x)) - \exp(iW_{h_2}(x))\right)\overline{\left(\exp(iW_{h_1}(x)) - \exp(iW_{h_2}(x))\right)}\mathrm{d}\Normal(x\mid 0,Q) \\
    & \hspace{0.5cm} \stp{=}{1} \int_\hs{H} 1 - \exp(iW_{h_1}(x)-iW_{h_2}(x)) - \exp(iW_{h_2}(x)-iW_{h_1}(x)) + 1\mathrm{d}\Normal(x\mid 0,Q) \\
    & \hspace{0.5cm} \stp{=}{2} 2 - \int_\hs{H} \exp(i \langle Q^{-\frac12}(h_1-h_2),x\rangle_\hs{H})\mathrm{d}\Normal(x\mid 0,Q) - \int_\hs{H} \exp(i \langle Q^{-\frac12}(h_2-h_1),x\rangle_\hs{H})\mathrm{d}\Normal(x\mid 0,Q) \\
    & \hspace{0.5cm} \stp{=}{3} 2 
        -  \exp(i \langle 0, h\rangle_\hs{H}) \exp\left(-\frac{1}{2}\langle Q \left(Q^{-\frac12}(h_1-h_2)\right), \left(Q^{-\frac12}(h_1-h_2)\right) \rangle_\hs{H}\right) \\
        & \hspace{1cm} - \exp(i \langle 0, h\rangle_\hs{H}) \exp\left(-\frac{1}{2}\langle Q \left(Q^{-\frac12}(h_2-h_1)\right), \left(Q^{-\frac12}(h_2-h_1)\right) \rangle_\hs{H}\right) \\
    & \hspace{0.5cm} \stp{=}{4} 2 - \exp\left(-\frac12\|h_1-h_2\|_{\hs{H}}^2\right)
        - \exp\left(-\frac12\|h_2-h_1\|_{\hs{H}}^2\right) \\
    & \hspace{0.5cm} = 2 - 2\exp\left(-\frac12\|h_1-h_2\|_{\hs{H}}^2\right),
\end{align*}
In \stpx{1}, we used the usual rules $e^{ix}\overline{e^{ix}}=|e^{ix}|^2=1$
and $e^xe^y=e^{x+y}$,
and in \stpx{2} we used the linearity of the integral, the fact that $\Normal(0,Q)$ is a probability distribution, and the definition of $h \mapsto W_h$ (recall that $h_1,h_2\in Q^{\frac12}(\hs{H})$).
For \stpx{3} we used again the characterization of the Gaussian distribution via the Fourier transform from Proposition \ref{prop:gaussian:existenceUniquenessFourier},
and in \stpx{4} we used
\begin{equation*}
    \langle Q(Q^{-\frac12}h),(Q^{-\frac12}h)\rangle_\hs{H}
    = \langle Q^\frac12Q^\frac12(Q^{-\frac12}h),(Q^{-\frac12}h)\rangle_\hs{H}
    = \langle Q^\frac12(Q^{-\frac12}h), Q^\frac12(Q^{-\frac12}h) \rangle_\hs{H}
    = \|h\|_\hs{H}^2,
\end{equation*}
which holds for all $h\in Q^{\frac12}(\hs{H})$.
Observe now that for $h_2\rightarrow h_1$, the righthand side of the above equality chain converges to zero, showing the continuity of the map 
$Q^{\frac12}(\hs{H})\ni h \mapsto \exp(i W_h) \in L^2(\hs{H},\Normal(0,Q);\C)$,
and a fortiori of all the maps
$Q^{\frac12}(\hs{H})\ni h \mapsto \exp(i \lambda W_h) \in L^2(\hs{H},\Normal(0,Q);\C)$
with $\lambda\in\R$.
The result now follows due to denseness of $Q^\frac12(\hs{H})$ in $\hs{H}$.
\end{proof}

\end{document}
